\let\negmedspace\undefined
\let\negthickspace\undefined
\documentclass[journal]{IEEEtran}
\usepackage[a5paper, margin=10mm, onecolumn]{geometry}
%\usepackage{lmodern} 
\usepackage{tfrupee} 

\setlength{\headheight}{1cm} 
\setlength{\headsep}{0mm}     

\usepackage{gvv-book}
\usepackage{gvv}
\usepackage{cite}
\usepackage{amsmath,amssymb,amsfonts,amsthm}
\usepackage{algorithmic}
\usepackage{graphicx}
\usepackage{textcomp}
\usepackage{xcolor}
\usepackage{txfonts}
\usepackage{listings}
\usepackage{enumitem}
\usepackage{mathtools}
\usepackage{gensymb}
\usepackage{comment}
\usepackage[breaklinks=true]{hyperref}
\usepackage{tkz-euclide} 
\usepackage{listings}                                        
\def\inputGnumericTable{}                                 
\usepackage[latin1]{inputenc}                                
\usepackage{color}                                            
\usepackage{array}                                            
\usepackage{longtable}                                       
\usepackage{calc}                                             
\usepackage{multirow}                                         
\usepackage{hhline}                                           
\usepackage{ifthen}                                           
\usepackage{lscape}

\begin{document}

\bibliographystyle{IEEEtran}
\vspace{3cm}

\title{8.2.37}
\author{AI25BTECH11003 - Bhavesh Gaikwad}
{\let\newpage\relax\maketitle}

\renewcommand{\thefigure}{\theenumi}
\renewcommand{\thetable}{\theenumi}
\setlength{\intextsep}{10pt} 


\numberwithin{equation}{enumi}
\numberwithin{figure}{enumi}
\renewcommand{\thetable}{\theenumi}


\textbf{Question}: Find the equation of the conic, that satisfies the given conditions.\\
vertex (-3, 0), directrix x + 5 = 0.

\textbf{Solution:}\\
Given conditions:\\
Vertex: $(-3, 0)$\\
Directrix: $x + 5 = 0$\\

$\Rightarrow$ Since we have a vertex and directrix, this conic is a parabola with eccentricity $e = 1$.\\

Using the formulation for a conic with directrix $\vec{n}^\top\vec{x} = c$, eccentricity $e$, and focus $\vec{F}$:\\

Given Parameters:

Normal vector: $\vec{n} = \myvec{1\\0}$

Directrix constant: $c = -5$

Eccentricity: $e = 1$

Focus: $\vec{F} = \myvec{-1\\0}$ [derived from equidistance property]\\

Calculating matrix $\vec{V}$
\begin{equation}
\vec{V} = \norm{\vec{n}}^2\vec{I} - e^2(\vec{n}\vec{n}^\top) = \vec{I} -\myvec{1\\0}\myvec{1 & 0} = \myvec{0 & 0\\0 & 1}
\end{equation}

Calculating $\vec{u}$
\begin{equation}
\vec{u} = ce^2\vec{n} - \norm{\vec{n}}^2\vec{F} = (-5)(1)\myvec{1\\0} - 1\myvec{-1\\0} = \myvec{-4\\0}
\end{equation}

Calculating f:
\begin{equation}
f = \norm{\vec{n}}^2\norm{\vec{F}}^2 - c^2e^2 = 1 - 25 = -24
\end{equation}

\newpage

\begin{equation}
\text{The eigenvalues of } \vec{V} \text{ is } \lambda_1 = 0
\end{equation}

Since the eigenvalue is zero, we confirm this is a parabola $(e = 1)$.


The general conic equation is:
\begin{equation}
\vec{x}^\top\vec{V}\vec{x} + 2\vec{u}^\top\vec{x} + f = 0
\end{equation}

Substituting our values in Equation 0.5:
\begin{align}
\myvec{x & y}\myvec{0 & 0\\0 & 1}\myvec{x\\y} + 2\myvec{-4 & 0}\myvec{x\\y} + (-24) &= 0\\
y^2 - 8x - 24 &= 0
\end{align}\\

Therefore, the Equation of the Conic is:
\begin{equation}
\boxed{y^2 = 8(x + 3)}
\end{equation}


\begin{figure}
    \centering
    \includegraphics[width=\columnwidth]{figs/fig1.png}
    \caption{Parabola}
    \label{fig:figs/fig1.png}
\end{figure}

\end{document}  
